% Annex A: source PDF pages 35--37 / printed pages 23--25.
% Fidelity pass 04: page-local spacing and source-line membership are tuned
% against the supplied IEEE PDF while retaining live, semantic LaTeX text.

\newpage
\begingroup
\addtolength{\topskip}{1.4pt}
% Source page 35 title stack.  The shared annex macros were intentionally kept
% unchanged because Annex B already matches its source page exactly.
\annexdesignation{A}
\vspace{-6.9pt}
\annexstatus{informative}
\vspace{-7.4pt}
\annextitle{Rationale}
\vspace{-8.0pt}
\annexsection{A.1 Federation rules}
\vspace{-2.3pt}

\annexsubsection{A.1.1 Rationale for Rule 1}

\textit{Federations shall have an HLA FOM, documented in accordance with the HLA OMT.}

\vspace{6.1pt}
The HLA is domain-independent and can be used to support federations for a wide variety of uses. The\linebreak
formalization of agreements for information exchange is an important element of the HLA. By formalizing\linebreak
the development of these agreements and requiring that the results be documented in a common format, the\linebreak
HLA provides the means for understanding the key elements of a federation and for assisting in the reuse of\linebreak
the federation, in whole or in part. Such reuse is a goal of the HLA. In addition, the FOM provides the basis\linebreak
for some of the data used to initialize the RTI for the federation.

\vspace{4.5pt}
\annexsubsection{A.1.2 Rationale for Rule 2}

\textit{In a federation, all simulation-associated object instance representation shall be in the federates, not in the\linebreak
RTI.}

\vspace{6.1pt}
One basic idea behind the HLA is to separate federate-specific functionality from general-purpose\linebreak
supporting infrastructure. A necessary part of this feature is that the supporting infrastructure have no\linebreak
information about the objects being simulated. If the opposite were allowed, and information about a given\linebreak
application were stored in the supporting infrastructure, that infrastructure would by definition be tailored\linebreak
for that specific application and no longer be general purpose. Federate functionality was separated from\linebreak
federation support services for several reasons. First, the RTI services are intended to be the basic set of\linebreak
broadly reusable capabilities needed to support federations across the widest range of users. These are\linebreak
essentially coordination and management services supporting federation operations, time coordination, data\linebreak
distribution, etc. Because they apply across a range of HLA applications, these services can be provided\linebreak
most cost effectively as services to the applications rather than as components of the applications. Second,\linebreak
separation of the RTI services has the added advantage of freeing the federates to focus on their primary\linebreak
objective of representing object instances to meet the needs of a user or application domain. This approach\linebreak
frees the developers of federates from investing time and resources in these basic common services.

\vspace{4.6pt}
\annexsubsection{A.1.3 Rationale for Rule 3}

\textit{During a federation execution, all exchange of FOM data among joined federates shall occur via the RTI.}

\vspace{6.1pt}
The reason for providing common RTI services to federations is to provide in common the needed basic\linebreak
functionality to permit coherency in data exchange among the joined federates. If a federation were to\linebreak
exchange data representing state changes of shared object instances or interactions outside of the RTI service\linebreak
suite, the coherency of the distributed application would be violated.

\endgroup
\newpage
\begingroup
\addtolength{\topskip}{-0.7pt}
\annexsubsection{A.1.4 Rationale for Rule 4}

\textit{During a federation execution, joined federates shall interact with the RTI in accordance with the HLA\linebreak
interface specification.}

\vspace{6.1pt}
By requiring a standardized, common interface between federates and the RTI, along with a common\linebreak
application programmer's interface (API), the HLA allows for independent development and\linebreak
implementation of federate applications. Federate developers can work independently and develop\linebreak
interfaces to the RTI without regard to RTI implementation, and RTI developments can proceed without\linebreak
explicit consideration of federate development. The separation of the interfaces from the requirements for\linebreak
federate data exchange allows for the reuse of a common interface specification across the broad spectrum\linebreak
of distributed M\&S applications, with specific application needs tailored through the FOM mechanism.

\vspace{4.5pt}
\annexsubsection{A.1.5 Rationale for Rule 5}

\textit{During a federation execution, an instance attribute shall be owned by at most one joined federate at any\linebreak
given time.}

\vspace{6.1pt}
The HLA also provides the mechanism during federation execution to dynamically transfer ownership of an\linebreak
instance attribute from one joined federate to another or for a joined federate to unconditionally divest its\linebreak
ownership of an instance attribute. In this second case, the instance attribute becomes ``unowned'' by all\linebreak
joined federates, hence the rule explicitly states that ``an instance attribute shall be owned by at most one\linebreak
joined federate at any given time.'' Joined federates can request ownership of ``unowned'' instance attributes.\linebreak
By defining ownership at the instance attribute level and providing the tools to hand off ownership during\linebreak
federation execution, the HLA provides a flexible tool set for using various combinations of joined federates\linebreak
to meet user needs.

\vspace{10.8pt}
\annexsection{A.2 Federate rules}

\vspace{-2.4pt}
\annexsubsection{A.2.1 Rationale for Rule 6}

\textit{Federates shall have an HLA SOM, documented in accordance with the HLA OMT.}

\vspace{6.1pt}
A major goal of the HLA is to support interoperability and reuse of simulations. The HLA provides this\linebreak
support by providing for reuse at the level of simulations (or, more generally, federates) and allowing access\linebreak
to the representations in those simulations.

\vspace{6.1pt}
Lack of cost-effective access to information about the object representations available in federates inhibits\linebreak
reuse. The requirement for a SOM addresses information access by requiring federates to document the\linebreak
minimum essential, salient aspects of their capabilities to allow for easy identification of the federates'\linebreak
potential application in new federations. Although the full set of information required by a potential user\linebreak
will go well beyond the SOM contents, providing easy access to characterizations of simulation based on\linebreak
reuse potential will enable users to decide more effectively whether to invest in further assessment of the\linebreak
simulations to their applications. There are certain circumstances, such as dynamically configured passive\linebreak
subscribers, where a federate's SOM is based on the current FOM.

\endgroup
\newpage
\begingroup
\addtolength{\topskip}{-1.05pt}
\annexsubsection{A.2.2 Rationale for Rule 7}

\textit{Federates shall be able to update and/or reflect any instance attributes and send and/or receive interactions,\linebreak
as specified in their SOMs.}

\vspace{6.1pt}
By designing federates from the outset with the ability to present internal objects/attributes/interactions as\linebreak
public, the mechanisms for reuse of the federate will be in place from the start.

\vspace{4.5pt}
\annexsubsection{A.2.3 Rationale for Rule 8}

\textit{Federates shall be able to transfer and/or accept ownership of instance attributes dynamically during a\linebreak
federation execution, as specified in their SOMs.}

\vspace{6.1pt}
By building in the functionality to transfer and accept ownership of instance attributes, federates designed in\linebreak
accordance with the HLA provide some of the basic structural tools to become federates in the widest\linebreak
possible range of future federations. With this capability, it is possible to allow a federate designed for one\linebreak
purpose to be coupled with one designed for another purpose to meet a new requirement.

\vspace{4.5pt}
\annexsubsection{A.2.4 Rationale for Rule 9}

\textit{Federates shall be able to vary the conditions (e.g., thresholds) under which they provide updates of\linebreak
instance attributes, as specified in their SOMs.}

\vspace{6.1pt}
By designing and building in the capability to vary the conditions under which they can export instance\linebreak
attributes, federates designed in accordance with the HLA provide some of the basic structural tools to\linebreak
become federates in the widest possibly range of future federations.

\vspace{4.5pt}
\annexsubsection{A.2.5 Rationale for Rule 10}

\textit{Federates shall be able to manage local time in a way that will allow them to coordinate data exchange with\linebreak
other members of a federation.}

\vspace{6.1pt}
The HLA time management structure is intended to support interoperability among federates that may rely\linebreak
on different internal concepts of time. For example, some federates might operate without any explicit\linebreak
notion of time or passage of time; other federates may need to run faster than real time; and some federates,\linebreak
as is the case with some discrete event simulations, might rely upon a scrupulously coordinated causal order\linebreak
of events throughout the federation.

\vspace{6.1pt}
The HLA provides a single, unifying approach to time management, in that it offers a ``toolbox'' of services\linebreak
that support (not guarantee) federations in managing their event synchronization in a way that will meet both\linebreak
the needs of the federation and the range of internal time management strategies of the member federates.

\vspace{6.1pt}
Different categories of federates typically use a subset of an RTI's services to implement their mechanisms.\linebreak
Federates do not explicitly indicate to the RTI how time flow is being managed (e.g., time stepped, event\linebreak
driven, external time synchronization, independent advance). Federates for which passage of time is implicit\linebreak
may not need to use the services described in Clause 8 of IEEE Std 1516.1-2010. However, a federate must\linebreak
ensure that its internal mechanisms permit coherent operation of the time management paradigm (or\linebreak
paradigms) used in the federation. Some of the considerations that occur when a federation combines\linebreak
multiple, differing time management paradigms are discussed in 8.1 of IEEE Std 1516.1-2010.

\endgroup
